Put \(p=4+\delta_0>4\), and fix a block \(b\), an arm \(a\), and an arbitrary pre-block history \(\mathcal F_{n,b-1}\). By \(\text{def:admissible-arrays}\), \(\text{ass:predictable-adjustment}\), and \(\text{ass:complete-block-randomization}\), the potential outcomes, adjustment values, residual lists, and quotas are fixed at this conditioning event, and the labels assigned to arm \(a\) form a simple random sample without replacement of size \(m=m_{nba}\) from \(N=n_{nb}\) labels. By \(\text{ass:consistency-no-interference}\), the observed residual on such a label is its fixed residual. Write
\[
X_i=R_{nbi}(a)-\bar R_{nba},\qquad
\mu=\frac1N\sum_{i=1}^N\delta_{X_i}=\mu_{nba},qquad
\nu=\frac1m\sum_{i:Z_{nbi}=a}\delta_{X_i}.
\tag{1}
\]
Jensen's inequality and \(\text{ass:residual-moments}\) give the history-uniform centered moment bound
\[
H:=\frac1N\sum_i|X_i|^p
\le2^{p-1}\left\{\frac1N\sum_i|R_{nbi}(a)|^p+|\bar R_{nba}|^p\right\}
\le2^pC_R.
\tag{2}
\]

Let \(C_T(x)=(-T)\vee(x\wedge T)\) for \(T\ge1\), and write \(\lambda^T=C_{T\#}\lambda\) for the pushforward of a measure \(\lambda\). We first control the bounded central part. If \(F_T\) and \(\widehat F_T\) are the distribution functions of \(\mu^T\) and \(\nu^T\), the exact first- and second-order inclusion probabilities of simple random sampling give, for every \(t\),
\[
\operatorname{Var}(\widehat F_T(t)\mid\mathcal F_{n,b-1})
=\frac{N-m}{m(N-1)}F_T(t)\{1-F_T(t)\}\le\frac1{4m}.
\tag{3}
\]
For finitely supported laws on the line, increasing matching gives
\(W_1(\nu^T,\mu^T)=\int|\widehat F_T(t)-F_T(t)|dt\): between consecutive ordered atoms, every plan must transport at least the cumulative mass imbalance across that cut, and increasing matching transports exactly that amount. Because both clipped laws are supported on \([-T,T]\), conditional Cauchy--Schwarz and (3) yield
\[
E\{W_1(\nu^T,\mu^T)\mid\mathcal F_{n,b-1}\}\le Tm^{-1/2}.
\tag{4}
\]

Let \(\bar X_S=m^{-1}\sum_{i:Z_{nbi}=a}X_i\). The exact finite-population sampling-variance identity and (2) imply
\[
E(\bar X_S^2\mid\mathcal F_{n,b-1})
=\left(\frac1m-\frac1N\right)S_{nba}^2\le C m^{-1}.
\tag{5}
\]
The observable measure is \(\widehat\mu_{nba}=T_{-\bar X_S\,\#}\nu\). Since clipping is one-Lipschitz,
\[
W_1(\widehat\mu_{nba}^T,\nu^T)\le W_1(\widehat\mu_{nba},\nu)=|\bar X_S|.
\]
Thus (4)--(5) and the triangle inequality give
\[
E\{W_1(\widehat\mu_{nba}^T,\mu_{nba}^T)\mid\mathcal F_{n,b-1}\}
\le CTm_{nba}^{-1/2}.
\tag{6}
\]
All constants in this proof may depend only on \(K,\varepsilon_0,p,C_R\).

We next prove the tail estimate used below, including observable recentering. From overlap, \(m/N=p_{nba}\ge\varepsilon_0\), so every realized sample satisfies
\[
\int|x|^p,d\nu(x)=\frac1m\sum_{i:Z_{nbi}=a}|X_i|^p
\le\frac NmH\le\varepsilon_0^{-1}H.
\tag{7}
\]
Jensen gives \(|\bar X_S|^p\le\int|x|^p,d\nu(x)\); hence
\[
\int|y|^p,d\widehat\mu_{nba}(y)
\le2^{p-1}\left\{\int|x|^p,d\nu(x)+|\bar X_S|^p\right\}
\le2^p\varepsilon_0^{-1}H.
\tag{8}
\]
Thus both the oracle and observable margins have a deterministic, history-uniform \(p\)-moment bound.

For later use, consider any \(K\) margins \((\lambda_a)_a\) satisfying \(\int|x|^p d\lambda_a\le H_0\), and let \(J^\pm(\lambda_1,\ldots,\lambda_K)\) denote the extrema of \(E(c^\top X)^2\) over their couplings. For any such coupling and \(X_a^T=C_T(X_a)\), expansion of the quadratic gives
\[
\begin{aligned}
&|(c^\top X)^2-(c^\top X^T)^2|\\
&\quad\le\sum_{a,d}|c_ac_d|\left\{
|X_a|\mathbf1(|X_a|>T)|X_d|+|X_a|,|X_d|\mathbf1(|X_d|>T)\right\}.
\end{aligned}
\tag{9}
\]
Hölder's inequality with exponents \(p,p,p/(p-2)\), followed by Markov's inequality, gives
\[
E\{|X_aX_d|\mathbf1(|X_a|>T)\}
\le H_0^{1/p}H_0^{1/p}\Pr(|X_a|>T)^{(p-2)/p}
\le H_0T^{-(p-2)}.
\tag{10}
\]
Contrast normalization implies \(\sum_a|c_a|\le\sqrt K\), so (9)--(10) show that clipping changes the cost of every coupling by at most \(2KH_0T^{-(p-2)}\). Clipping an original coupling gives a coupling of the clipped margins. Conversely, because these measures are finite, a coupling of the clipped margins can be lifted by drawing each original coordinate from its conditional empirical law given its clipped value. Applying the same bound to both directions and to minimizers and maximizers proves
\[
|J^\pm(\lambda_1,\ldots,\lambda_K)-J^\pm(\lambda_1^T,\ldots,\lambda_K^T)|
\le C H_0T^{-(p-2)}.
\tag{11}
\]
This is the required tail error; no optimizer uniqueness is used.

For margins supported on \([-T,T]\), the quadratic objective is Lipschitz in coordinatewise \(W_1\). Indeed, for \(x,y\in[-T,T]^K\),
\[
|(c^\top x)^2-(c^\top y)^2|
\le2T\|c\|_1\sum_a|c_a||x_a-y_a|.
\tag{12}
\]
Starting from any coupling of one set of margins, glue to each coordinate an optimal \(W_1\) coupling to the corresponding second margin. Equation (12), followed by the same construction in reverse, bounds the difference of both extrema by \(C_KT\sum_aW_1(\lambda_a,\widetilde\lambda_a)\). Apply this to the clipped oracle and observable margins, take conditional expectations, and use (6). Then use (2), (8), and (11) on both sides of this central comparison. For either sign,
\[
E\{|\widehat J_{nb}^{\pm}-J_{nb}^{\pm}|\mid\mathcal F_{n,b-1}\}
\le C\left\{T^2q_n^{-1/2}+T^{-(p-2)}\right\}.
\tag{13}
\]
Here count growth supplied \(m_{nba}\ge q_n\), and the fixed number of arms was absorbed into \(C\). Equation (13) explicitly separates the bounded central error from the moment-tail error.

The diagonal variance component is faster. Put
\[
A_S=\frac1m\sum_{i:Z_{nbi}=a}X_i^2,qquad A=\frac1N\sum_iX_i^2.
\]
Applying the same exact sampling-variance identity to the finite list \((X_i^2)_i\), and using the fourth-moment consequence of (2), gives
\[
E(|A_S-A|\mid\mathcal F_{n,b-1})\le Cm^{-1/2}.
\tag{14}
\]
For \(m,N\ge2\),
\[
\widehat S_{nba}^2=\frac m{m-1}(A_S-\bar X_S^2),
\qquad S_{nba}^2=\frac N{N-1}A.
\]
Equations (2), (5), (14), and the elementary bounds on the two finite-sample factors therefore imply
\[
E\{|\widehat S_{nba}^2-S_{nba}^2|\mid\mathcal F_{n,b-1}\}
\le Cm_{nba}^{-1/2}.
\tag{15}
\]

Now take
\[
T=q_n^{1/(2p)},qquad
\kappa_\star=\frac{p-2}{2p}=\frac{2+\delta_0}{2(4+\delta_0)}.
\tag{16}
\]
For \(q_n\ge2\), this choice is admissible and balances the two terms in (13):
\[
T^2q_n^{-1/2}=T^{-(p-2)}=q_n^{-\kappa_\star}.
\tag{17}
\]
Let \(t_{nb}=w_{nb}^2/n_{nb}\), so \(\sum_bt_{nb}=s_n^2\). From overlap,
\(m_{nba}^{-1}\le(\varepsilon_0n_{nb})^{-1}\), and for \(n_{nb}\ge2\),
\((n_{nb}-1)^{-1}\le2n_{nb}^{-1}\). Combining (13), (15), (17), \(\sum_ac_a^2=1\), and the two envelope definitions gives
\[
E\left[|\widehat L_{nb}-L_{nb}|+|\widehat U_{nb}-U_{nb}|
\mid\mathcal F_{n,b-1}\right]
\le Ct_{nb}\{q_n^{-1/2}+q_n^{-\kappa_\star}\}
\le C't_{nb}q_n^{-\kappa_\star},
\tag{18}
\]
where the four symbols with subscript \(nb\) denote the corresponding block summands.

The bound (18) holds with the same constant at every possible pre-block history. Sum over blocks and apply the tower property under the original sequential randomization law; no maximum or union bound over blocks or histories appears. Since \(\sum_bt_{nb}=s_n^2\),
\[
\sup_{\mathbb D\in\mathfrak D}E_{\mathbb D}
\frac{|\widehat L_n-L_n|+|\widehat U_n-U_n|}{s_n^2}
\le C_{\mathrm{ep}}q_n^{-\kappa_\star}
=C_{\mathrm{ep}}r_n
\quad(q_n\ge2).
\tag{19}
\]
If \(q_n<2\), the loss in the theorem is truncated at one and
\(1\le2^{\kappa_\star}q_n^{-\kappa_\star}\); increasing \(C_{\mathrm{ep}}\) covers these total-convention cases. This proves the expectation bound for every \(n\).

For every nonnegative \(X\) and \(\eta>0\),
\(\mathbf1\{X>\eta\}\le(1\wedge X)/(1\wedge\eta)\). Applying this inequality to (19) proves the displayed probability bound. The Hausdorff distance between real intervals is
\(\max\{|\widehat L_n-L_n|,|\widehat U_n-U_n|\}\), which is bounded by the displayed sum, so it obeys the stated bound with the same constant. Finally, \(q_n\to\infty\) and \(\kappa_\star>0\) imply \(r_n\to0\). Every conditioning step above is at \(\mathcal F_{n,b-1}\); none freezes \(\mathcal F_{n,B_n}\).